\documentclass{article}
\usepackage{amsmath,amssymb,amsthm,hyperref}
\newtheorem{lemma}{Lemma}
\begin{document}
\title{Vector memories, fixed readout spans, and the remaining numerator cost}
\maketitle

\paragraph{Status and scope.}
The feature-span statements below are kernel checked in
\texttt{JointCellVectorMemoryCapacity.lean}. The rational identities and
capacity discussion are written proofs or explicitly identified open
steps. They do not improve the asymptotic upper bound for the original
worst-case head complexity and do not change its definition.

\paragraph{A positive vector-memory identity.}
Let $u_1,\ldots,u_r$ be distinct real roots below every pole $p$ and set
$a_l(u)=\prod_{k\ne l}(u_l-u_k)^{-1}$. Then
\[
 \frac1{\prod_l(p-u_l)}=\sum_l\frac{a_l(u)}{p-u_l},
 \qquad
 -\frac{\partial}{\partial p}\frac1{\prod_l(p-u_l)}
   =\sum_l\frac{a_l(u)}{(p-u_l)^2}.
\]
The first equality is partial fractions; differentiation gives the second.
The roots can be chosen independently for each table column, in any fixed
compact region separated from the poles. All displayed denominators are
positive there. This separation alone is not whole-cube admissibility;
the affine extension of each pole and root must retain a positive gap
on every Boolean input.

The first identity is not automatically an $r$-head compiler when the
original numerator is a row-dependent amplitude $A(x)$: its residue
$A(x)a_l(u_j)$ is a row-column product. The squared-kernel identity has
a two-head tangent implementation for each term, because
\[
 \frac{A(x)}{\varepsilon}
 \left(\frac1{D(x)-u_{lj}-\varepsilon a_l(u_j)}
             -\frac1{D(x)-u_{lj}}\right)
 \longrightarrow
 \frac{A(x)a_l(u_j)}{(D(x)-u_{lj})^2}.
\]
On a finite cube with an existing positive gap, sufficiently small
$\varepsilon$ preserves positivity. This identity still requires both
$A$ and $D$ to be affine in the actual input encoding. A tensor numerator
is not made affine by writing it as a symbol $A$.

\begin{lemma}[Fixed readout span]
Let $X$ be any set, let $\phi:X\to\mathbb R^d$ be any map, and fix
$v_1,\ldots,v_a\in\mathbb R^d$. If the vectors
$(v_b\cdot\phi(x))_{b=1}^a$, as $x$ varies, realize every strict sign
word in $\{-1,1\}^a$, then $a\le d$.
\end{lemma}
\begin{proof}
If $a>d$, there is a nonzero vector $z$ such that $\sum_bz_bv_b=0$.
Choose the strict sign word agreeing with $z_b$ at every nonzero
coordinate. Its proposed realization has
$\sum_b z_b(v_b\cdot\phi(x))>0$, contrary to the linear relation.
Zero coordinates of $z$ may be assigned either sign.
\end{proof}
The exact Lean declaration is
\texttt{jointCell\_vector\_memory\_feature\_capacity}.
The memory $x$ may be a tuple of roots, a vector of poles, or any other
nonlinear or label-adaptive object. The lemma only fixes the readout span.

\paragraph{Additive tensor residues.}
For $\alpha\in\{1,\ldots,q\}^c$, suppose
$\eta_b(\alpha)=\sum_l\eta_{bl}(\alpha_l)$ and the kernel
$K_\alpha(u)$ is independent of the query bit $b$. Then
\[
 \sum_\alpha\eta_b(\alpha)K_\alpha(u)
 =\sum_{l,z}\eta_{bl}(z)
      \sum_{\alpha:\alpha_l=z}K_\alpha(u).
\]
Thus $a\le cq$ is necessary for all strict words. Removing redundant
constant coordinates improves this to $a\le1+c(q-1)$; the checked
declaration \texttt{jointCell\_additive\_vector\_memory\_capacity}
uses the simpler $cq$ bound. An additional independently chosen row bias
adds at most one feature. No scalar zero-count theorem is needed.

There is a second encoding charge. Separate one-hot coordinates for the
bit $b$ and the $c$ plane points do not make $\eta_{bl}(\alpha_l)$
affine. If only the first plane block jointly encodes $(b,p_1)$, the
admissible numerator instead has the form
\[
 A_b(\alpha_1)+\sum_{l>1}A_l(\alpha_l).
\]
The common part contributes one feature, and the bit-dependent part
contributes $q$ features. Hence $a\le q+1$, as checked by
\texttt{jointCell\_one\_bit\_block\_vector\_memory\_capacity}.
Replicating $(b,p_l)$ in every block uses $acq^2$ physical coordinates
and creates a further consistency-cover problem: the intended tuples
with equal bit labels form only an $a^{1-c}$ fraction of the full product
of the encoded blocks. No cover saving is asserted for that restriction.

More generally, a factorization
$\eta_b(\alpha)=\sum_{t=1}^R v_{bt}w_t(\alpha)$ with a bit-independent
kernel requires $a\le R$, up to explicitly included bias features.
Consequently a small rank in the bit coordinate is not a free repair of
the numerator condition. Bit-dependent pole shifts or kernels escape
this lemma and require a new word-capacity argument.

\paragraph{What the scalar tests do and do not show.}
For rank-one residues on additive pole grids, the Laplace density is a
product of $c$ exponential sums with $q$ terms each and has at most
$c(q-1)$ sign changes. With all other roots fixed, dividing residues by
positive factors $(p-u_l)$ corresponds to repeated exponential tail
integration of that density; these operations do not increase sign
variation. The derivative product kernel inserts the positive factor
equal to total Laplace time before these integrations. Thus every
one-coordinate section has at most $c(q-1)$ zeros.

This fiber statement is not a multivariate arrangement bound. Even
monotone graphs can have many pairwise intersections. In particular, it
does not justify a bound of $O(c\log(cq))$ on the number of available
bits. Sums of two product densities connect to the two-product case of
the real tau problem; no polynomial zero bound is used here.

For $q=2$, additive residues instead give a density proportional to
$A+\sum_l\beta_l\tanh(w_lt/2)$ after division by a positive product.
The scale family $\tanh(wt)$ is not a Chebyshev system of every order.
With $w_i=1+i/1000$ and $t_j=T(1+j/1000)$ for $0\le i,j<10$, the
determinant of $[\tanh(w_it_j)]$ is negative at $T=1$ and positive at
$T=2$. The file \texttt{JointCellTanhChebyshevCertificate.py} verifies
both signs using rational intervals and a Taylor remainder bound for the
exponential. This is a finite computation, not a Lean declaration. By
continuity, a nonzero linear combination has ten distinct positive
zeros at an intermediate scale. This only refutes the sharp nine-zero
Chebyshev claim, not an $O(cq)$ bound or any polynomial bound.

\paragraph{Approximate lower-face filters, an open alternative.}
A degree of order $\min\{c,\sqrt{c\log(1/\epsilon)}\}$ can approximate
an endpoint indicator on the integer weights $0,\ldots,c$; see de Wolf,
\emph{A note on quantum algorithms and the minimal degree of
$\epsilon$-error polynomials for symmetric functions},
\url{https://arxiv.org/abs/0802.1816}.
Replacing exact degree-$c$
filters by such a polynomial would be useful only if a global compiler
tolerated an inverse-polynomial error $\epsilon$. Descending sign repair
does not presently supply this: earlier scores and later coefficients
may have arbitrarily large ratios. A simultaneous interaction argument
would need an explicit bound on the local score away from its target,
relative to its strict target margin, followed by a bound on the global
interaction operator. Neither estimate has been proved for a decoder
with a growing gain. The small average number of proper faces in a
random weighted cover does not by itself provide either estimate.

\paragraph{Numerical screen.}
\texttt{JointCellFiniteVectorCauchySearch.py} tests whether a reciprocal
rank-three matrix can be fitted by a small sum of positive Cauchy banks.
It uses nonlinear least squares and is only an exploratory screen.
A failed optimizer is not an impossibility certificate; a small residual
is not an exact decomposition or an admissible original-model compiler.
No theorem in this note depends on this search.
\end{document}
